% Keterbatasan penelitian
\section{Keterbatasan Penelitian}

Sampel dihimpun secara sukarela dari pengguna platform X sehingga condong ke arah skor tinggi pada kedelapan subtes (Tabel~\ref{tab:deskriptif}), dengan proporsi nilai di bawah median nasional yang minim. Parameter pembandingnya pun bersumber dari dokumen tidak resmi berupa materi internal Panitia SNPMB yang beredar di media sosial, sehingga ketidakakuratan pada data dasar itu berpotensi memengaruhi validitas komparasi parameter. Estimasi paling peka terhadap parameter ini (lihat Hasil), dan skenario pengali $SD$ yang diuji bersifat ilustratif sehingga tidak menjamin mencakup nilai sebenarnya. Dengan demikian, interval \textit{bootstrap} hanya mencerminkan sebagian ketidakpastian.

Audiens akun media sosial sumber data telah tersaring secara sosiologis sejak awal. Perluasan pengumpulan data melalui \textit{framing} yang lebih inklusif tidak mereduksi bias seleksi (lihat Pembahasan).

Data terbatas pada platform X dan pada satu tahun ujian (2026), sehingga karakteristik bias yang teramati belum tentu tergeneralisasi pada media atau komunitas daring lain. Pelaporan dua skor LBI juga merupakan format khas tahun itu. Oleh karena itu, temuan belum tentu berlaku bagi UTBK-SNBT tahun lain maupun ujian seleksi serupa di Indonesia.

Prosedur pembersihan (lihat Metode) telah mengeksklusi indikasi rekayasa kecerdasan buatan, skor bulat sempurna, dan anomali tata letak visual, tetapi peluang lolosnya entri palsu tidak dapat disingkirkan sepenuhnya.

Koreksi Case III bertumpu pada asumsi homogenitas korelasi. Formula Thorndike Case III mengasumsikan struktur korelasi yang relatif seragam antarpasangan subtes untuk menyederhanakan $SD$ skor komposit populasi menjadi satu nilai tunggal. Diagnostik deviasi $SD$ tersirat (Tabel~\ref{tab:diagnostik}) menunjukkan bahwa asumsi ini tidak terpenuhi secara seragam pada seluruh subtes, sehingga estimasi $\bar{R} \approx 0,425$ paling tepat ditafsirkan sebagai rerata agregat populasi. Koreksi ini juga mengasumsikan hubungan linear antara skor komposit dan tiap subtes dengan sebaran galat yang sama di seluruh rentang skor, dan asumsi itu dapat terlanggar apabila korelasi antartes berubah menurut tingkat kemampuan (lihat Pembahasan). Simulasi skenario SLODR menunjukkan bahwa dalam model satu faktor sederhana, koreksi Case III menaksir rendah $\bar{R}$ dan $S_Z$ bila korelasi menurun pada skor tinggi (Tabel~\ref{tab:slodr_case3}). Skenario itu bersifat ilustratif karena bentuk dan besaran kemiringannya tidak diestimasi dari data, dan data paper ini tidak memungkinkan pemisahan SLODR dari \textit{restriction of range}.

Tabel translasi persentil memakai asumsi normal sebagai model kerja, padahal statistik nasional tidak resmi menunjukkan distribusi miring ke kanan dengan ekor kiri terpotong. Analisis sensitivitas dengan marginal Johnson SB dan \textit{copula} Gaussian menunjukkan bahwa asumsi ini berdampak kecil di sekitar $\pm 1\,SD$ tetapi mengubah peluang dan skor di ekor secara berarti. Hasil analisis itu bersifat indikatif karena tujuh statistik per subtes tidak mengidentifikasi bentuk distribusi secara unik dan dependensi pada simulasi diasumsikan seragam.